\documentclass[letterpaper,11pt]{article}
\usepackage[a4paper,margin=0.45in,top=0.3in,bottom=0.25in]{geometry}
\usepackage{fontspec}
\IfFontExistsTF{Latin Modern Roman}{\setmainfont{Latin Modern Roman}[SmallCapsFont={Latin Modern Roman Caps}]}{\setmainfont{lmroman10}[Extension=.otf,UprightFont=*-regular,BoldFont=*-bold,ItalicFont=*-italic,BoldItalicFont=*-bolditalic,SmallCapsFont=lmromancaps10-regular]}% fallback by file name (tectonic)
\usepackage{titlesec,enumitem,tabularx,xcolor,graphicx}
\usepackage[hidelinks,pdfauthor={Firstname Lastname}]{hyperref}
\pagestyle{empty}
\setlength{\parindent}{0pt}
\raggedbottom\raggedright
\titleformat{\section}{\vspace{-6pt}\scshape\raggedright\large}{}{0em}{}[\titlerule\vspace{-5pt}]
\newcommand{\entry}[4]{\item\begin{tabular*}{0.97\textwidth}[t]{l@{\extracolsep{\fill}}r}\textbf{#1} & #2 \\ \textit{\small#3} & \textit{\small #4}\\\end{tabular*}\vspace{-7pt}}
\newcommand{\proj}[2]{\item\begin{tabular*}{0.97\textwidth}{l@{\extracolsep{\fill}}r}\small#1 & #2\\\end{tabular*}\vspace{-7pt}}
\newcommand{\bul}[1]{\item\small{#1\vspace{-2pt}}}
\newcommand{\entriesStart}{\begin{itemize}[leftmargin=0.15in,label={}]}
\newcommand{\entriesEnd}{\end{itemize}}
\newcommand{\bulStart}{\begin{itemize}[leftmargin=0.3in,label=\textbullet,itemsep=-1.5pt]}
\newcommand{\bulEnd}{\end{itemize}\vspace{-5pt}}
\newcommand{\header}[4]{%
\begin{minipage}[t]{0.09\textwidth}\mbox{}\end{minipage}\hfill
\begin{minipage}[t]{0.79\textwidth}\centering
{\Huge\scshape Firstname Lastname}\\[2pt]
\large #1\\[1pt]
\if\relax\detokenize{#2}\relax\else{\small\textbf{#2}}\\[1pt]\fi
\small #3\\
\small \href{https://linkedin.com/in/your-profile}{\underline{linkedin.com/in/your-profile}} $|$ \href{https://github.com/your-github}{\underline{github.com/your-github}}
\end{minipage}\hfill
\begin{minipage}[t]{0.10\textwidth}\raggedleft\vspace{0pt}\IfFileExists{#4}{\fbox{\includegraphics[width=\linewidth]{#4}}}{}\end{minipage}
\vspace{2pt}}

\begin{document}
% Deutsche Version: dieselben Fakten wie cv_en.tex und CLAUDE.md Abschnitt 9. Eine Seite.
\header{Ihre Berufsbezeichnung}{%%TARGET%%}{Stadt $|$ +00 000 000 00 00 $|$ \href{mailto:you@example.com}{\underline{you@example.com}}}{photo.jpg}
\section{Profil}
\small %%SUMMARY%%
\section{Berufserfahrung}
\entriesStart
\entry{Position}{2025 -- heute}{Firma oder ``Freiberuflich''}{Stadt oder remote}
\bulStart
\bul{Was Sie gebaut haben, für wen, mit einer \textbf{echten Zahl}}
\bul{Zweites Ergebnis mit den eingesetzten Werkzeugen}
\bulEnd
\entry{Frühere Position}{MM/JJJJ -- MM/JJJJ}{Firma}{Stadt, Land}
\bulStart
\bul{Eine Zeile zu Ihren Aufgaben}
\bulEnd
\entriesEnd
\section{Projekte}
\entriesStart
\proj{\textbf{Projektname} $|$ \emph{Python, Docker, GitHub Actions}}{\href{https://github.com/your-github/project}{\underline{github.com/your-github/project}}}
\bulStart
\bul{Was es macht und warum es zählt}
\bulEnd
\entriesEnd
\section{Kenntnisse}
\begin{itemize}[leftmargin=0.15in,label={}]\small\item{
\textbf{KI und LLM}: ...\\
\textbf{Automatisierung}: ...\\
\textbf{Sprachen}: ...\\
\textbf{Frameworks und Tools}: ...}
\end{itemize}
\section{Ausbildung}
\entriesStart
\entry{Schule, Hochschule oder Kurs}{JJJJ}{Abschluss oder Zertifikat}{Stadt oder online}
\entriesEnd
\section{Weiteres}
\begin{itemize}[leftmargin=0.15in,label={}]\small\item{
\textbf{Sprachen}: Sprache (Niveau), ...\\
\textbf{Verfügbarkeit}: remote oder vor Ort; Arbeitserlaubnis; Eintrittstermin}
\end{itemize}
\end{document}
